\documentclass[../principal.tex]{subfiles}
\graphicspath{{\subfix{../imagenes/}}}

\begin{document}

\ifSubfilesClassLoaded{
    \part{Máquinas}
}{}

\chapter{Amplificadores}

Típicamente las señales eléctricas que salen de un módulo de sonido necesitan ser amplificadas antes de ser alimentadas a un parlante para ser escuchadas.

Los componentes que hacen ese trabajo se llaman amplificadores, y se representan con la letra $A$.

Un amplificador tiene un parámetro llamado ganancia, que se representa con la letra $G$. La ganancia es un número adimensional que indica cuánto se amplifica la señal de entrada.

Los amplificadores no son perfectos, y tienen una ganancia máxima, un ancho de banda, y todo tipo de limitaciones, que dependerán del uso que queramos, el precio que queramos, o la estabilidad que queramos.

\section{VCA}

Existe en el rubro de máquinas sonoras un tipo de amplificador del tipo controlado por voltaje, que recibe el nombre VCA por sus siglas en inglés.

\newpage

\end{document}

